\clearpage
\section{The sharp inequality and strictness above the finite range}
\label{short:composition}

\begin{theorem}[The sharp iid Berry--Esseen inequality]
\label{short:global}
Let $X_1,X_2,\ldots$ be independent copies of a random variable with
$\E X_1=0$, $\E X_1^2=1$ and $\rho=\E|X_1|^3<\infty$.
For every integer $n\ge1$,
\[
 \sup_{H\in\R}\left|
 \Pp\!\left(\sum_{i=1}^nX_i\le H\right)-\Phi(H/\sqrt n)
 \right|\le\frac{\CE\rho}{\sqrt n},
 \qquad \CE=\frac{3+\sqrt{10}}{6\sqrt{2\pi}}.
\]
Equivalently, $C_n\le\CE$ for every $n\ge1$. Moreover,
$C_n<\CE$ for every $n\ge700$, and $\sup_{n\ge1}C_n=\CE$.
\end{theorem}

\begin{proof}[Proof of the upper bound]
The small-index comparisons and the finite induction establish
$C_j\le\CE$ for $1\le j\le42$; their complete numerical checks
and same-law input joins are recorded in
Appendix~\ref{revision:reproducibility}. Suppose that the upper bound
fails, and choose the least index $n$ with $C_n>\CE$.
Then $n\ge43$, both deleted predecessors are sharp, and
Corollary~\ref{ext:first-failure-entry} gives
\[
 0<n(C_n-\CE)<\frac{3003}{10250}<\frac3{10},\qquad
 C_n<\frac{73}{175}<\frac{4226}{10000},\qquad
 n(C_{n-1}-C_n)_+=0.
\]
Since $\CE<41/100<4253/10000$, both deleted-predecessor ceilings
also hold. These bounds are stronger than the hypotheses of the finite local exclusion
when $43\le n\le48$, and of the uniform local exclusion when
$n\ge49$. Propositions~\ref{ext:local42-input} and
\ref{ext:local47-input} exclude both alternatives. Thus the upper
bound holds at every index. This is the first-failure composition
of Theorem~\ref{ext:bridge-composition}.
\end{proof}

The weak upper bound already proves the conjectured inequality.
The compact fixed-target comparison supplies the stated strictness.
We give the interface explicitly because it uses a different
optimization problem from the finite induction.

\begin{theorem}[Exclusion of a nonnegative-excess maximizer]
\label{short:acquired-exclusion}
Let $n\ge700$ and suppose $C_{n-1}\le\CE$. No pair $(P,H)$
attaining the fixed-target maximum at $\CE$ has $a(P,H)\ge0$.
\end{theorem}
\begin{proof}
The structural inequality~\eqref{v9:entry:structural} and
Theorem~\ref{v9:entry:theorem} give $nD\le32.20665378$ and
$|H|\le1.83605680$ for the same compactly supported law.
Propositions~\ref{short:acq:geometry} and \ref{short:acq:strip}
then give a central gap and the bounds below. In
\eqref{revision:canonical-decomposition}, identify the local theorem's
$B$ with $B_0=\ind_{\{X>0\}}$. The variable
$(B+U-p)/\sqrt{pq+s}$ is exactly the original $X$, and $H$ remains
its raw sum threshold.
\begin{center}\small
\begin{tabular}{@{}P{.49\linewidth}P{.44\linewidth}@{}}
\toprule
Bound just obtained & Hypothesis of Theorem~\ref{short:count:local} \\
\midrule
$3/8<p<93/200$ & $p\in[3/8,93/200]$ \\
$a_0,a_1\ge0$, $a_0+a_1<1/250$ & Nonnegative inward endpoints, sum at most $1/250$ \\
$-3/5<U_0\le a_0$, $-a_1\le U_1<3/5$ & The stated closed conditional-support bounds \\
$0\le\lambda<7/20$ & Positive-noise interval or separate zero-noise case \\
$|H|\le1.83605680<46/25$ & The full threshold interval \\
Independent copies of $X$ & The iid sum in the local comparison \\
\bottomrule
\end{tabular}
\end{center}
Positive noise gives $a(P,H)<-\lambda/100$. At zero noise the law
is Bernoulli, so Schulz's strict finite-binomial bound gives
$a(P,H)<0$. Both contradict $a(P,H)\ge0$.
\end{proof}

\begin{proof}[Proof of strictness and sharpness in Theorem~\ref{short:global}]
The proved upper bound supplies $C_{n-1}\le\CE$ for every $n\ge700$.
If $C_n\ge\CE$, fixed-target attainment and reflection give a
maximizer with $a\ge0$, contrary to the preceding theorem.
Hence $C_n<\CE$ for $n\ge700$.
Finally, Esseen's fixed skew-Bernoulli limit
\eqref{ms:bernoulli-limit} gives $\sup_n C_n\ge\CE$.
Together with the upper bound, this proves equality for the global
supremum.
\end{proof}

The two local comparisons have different purposes. The first-failure
comparison proves the weak bound at every index; the compact comparison
excludes equality from $700$ onward. Neither argument assumes that
an unrestricted extremizer is a two-point law. The certificates verify
inequalities on whole parameter regions for the original marginal,
after the analytic acquisition of those regions.
